\scanpage{15}
\deleted{$k=$ corps premier}

\deleted{$K$ corps}\quad\deleted{$k=$ corps premier}, $k=\unsure{\mathbb Q|Q}$

$\widehat{\unsure{\mathscr X|\mathcal X|\mathfrak X}}$ droite projective sur $\bar K$
\[
S\subset\mathscr X(K)\qquad S\ \text{partie finie}
\]
\[
\mathscr X=\widehat{\mathscr X}\setminus S
\]
$\mathscr X'$ rev. abélien [``modéré''] maximal de $\mathscr X$.

A) Si $\mathrm{Card}\,S\in\{0,1\}$, alors $\mathscr X'=\mathscr X$ ...

B) Si $\mathrm{Card}\,S=2$, alors $\mathscr X$ est \inserted{(\unsure{canoniqt|canonique})} un torseur sous une forme de $\mathrm{Gm}_{\overline{\mathbb Q}}$ ; \deleted{\unsure{et}} choisir un des pts de $S$ comme $\infty$ \inserted{(l’autre étant pris comme $0$)} revient à choisir un isom du \deleted{\unsure{ce}} groupe d’op. avec $\mathrm{Gm}$ ; \deleted{\unsure{et|choisir}} choisir une ``origine'' (sect $1$) de $\mathscr X$ revient à trivialiser ce torseur. \deleted{Ceci étant fait, on \unsure{trouve}} \inserted{En $\overline{\mathrm m}$ temps} on a une descente de la situation $(\widehat{\mathscr X},\mathscr X)$ de $\overline{\mathbb Q}$ à $\mathbb Q$, en $(\widehat X,X)$
\[
\widehat X=\mathbb P^1_{\mathbb Q},\quad X=\mathbb P^1_{\mathbb Q}\setminus\{0,\infty\}=\mathrm{Gm}_{\mathbb Q}.
\]
Ici
\[
\unsure{\mathcal G|\mathscr G}=\Gamma\mathbin{\underset{1/2}{\cdot}}\overline{\mathbb Q}^{*}
\]
$\Gamma$ opérant sur $\overline{\mathbb Q}^{*}$ de la façon évidente. Le rev. \unsure{kummérien|kummerien}\inserted{-type} d’indice $n$ de $X$ est $X_n^{=X}$, avec $\deleted{\unsure{X}}X_n\xrightarrow{t\mapsto t^n}X$, rev principal de groupe $\mu_{n\,\mathbb Q}=\mu_n$. On a, posant
\diagram{\begin{tikzcd}[ampersand replacement=\&,row sep=tiny]
X_n\arrow[d,"f_n(t)=t^n"]\\X
\end{tikzcd}\qquad
\begin{tikzcd}[ampersand replacement=\&,row sep=tiny]
\mathcal G_n=\mathrm{Aut}_{\unsure{K|k}}(\mathscr X_n\mathbin{-}\mathscr X)\arrow[d,"\unsure{\wr|2}"]\\
\mathrm{Aut}_{k}(\mathscr X_n)
\end{tikzcd}}
\[
\mathcal G_n^{\unsure{\prime}}=\Gamma\mathbin{\underset{1/2}{\cdot}}\overline{\mathbb Q}^{*}.
\]
L’homomorphisme canonique
\diagram{\begin{tikzcd}[ampersand replacement=\&]\mathcal G_n\arrow[r]\&\mathcal G\end{tikzcd}}
étant l’identité sur $\deleted{\unsure{\mathcal G}}\Gamma$, et l’élévation à la puiss. \unsure{n:ième|nième|n-ième} sur $\unsure{\mathcal G|G|\mathfrak G}$
% archive sheet number: 14
